% The appendices of the notes.  This deck made no searches of its own, so
% there is only Bilaga A: the method, written for the students, with a
% table saying where each of the material's research claims is checked,
% and a closing section listing what is left undone.  A claim chapter
% (which would carry the project's \chapterprecis sentence) is added when
% a search has actually been carried out here.

\chapter{Att kontrollera ett påstående}
\label{app:verifiera}

Den här övningen gör, som allt undervisningsmaterial, en rad påståenden.
De flesta går att pröva vid tangentbordet: skriv av programmet, kör det
och se efter.  Det är faktiskt hela poängen med
\cref{sec:fem-fel,sec:anrop}.  Men några påståenden går inte att köra ---
de handlar om hur man bäst lär ut något, eller om hur nybörjare faktiskt
tänker --- och sådana måste kontrolleras mot litteraturen.

Den här övningen har inte gjort någon egen sökning.  Varje påstående den
gör är antingen redan kontrollerat i en av veckans föreläsningar, och
pekas då ut där, eller ännu inte kontrollerat, och räknas då upp i
\cref{sec:fortsatt}.  \Cref{tab:paastaenden} visar vilket som är vilket.
Bilagan i övrigt beskriver arbetsgången, som är densamma oavsett vem som
går den: du kommer att behöva den varje gång du läser en lärobok, en
webbsida eller ett svar från en språkmodell, och varje gång du själv
skriver en rapport.

\begin{table}[htbp]
  \centering
  \caption{Övningens påståenden om forskning och om lärande, och var
    vart och ett av dem kontrolleras.  Inget av dem kontrolleras i den
    här bilagan: övningen har inte gjort någon egen sökning.}
  \label{tab:paastaenden}
  \begin{tabular}{@{}p{0.50\linewidth}p{0.44\linewidth}@{}}
    \toprule
    Påstående & Kontrolleras \\
    \midrule
    Att det lönar sig att försöka själv innan man läser
      lösningsförslaget, och att vinsten kommer av jämförelsen och inte
      av väntan
      & Redan belagt i föreläsningen \emph{Upprepningar}, bilaga~C \\
    Att kontrast måste upplevas samtidigt för att en aspekt ska bli
      urskiljbar, och att en öppen fråga visar vilka aspekter studenten
      redan urskiljer
      & Andrahandskälla, se kommentaren efter tabellen \\
    Att det är ett katalogiserat fel att glömma ändra räknaren i en
      \mintinline{python}{while}-slinga
      & Redan belagt i föreläsningen \emph{Upprepningar}, bilaga~B \\
    Att ett program inte ska upprepa sig --- DRY
      & Redan belagt i föreläsningen \emph{Algoritmiskt tänkande},
        bilaga~E \\
    Att en funktion eller en modul bör ha ett enda ansvar --- SRP
      & Redan belagt i föreläsningen \emph{Funktioner}, bilaga~B \\
    Fem råd och en missuppfattning som övningen ger utan belägg
      & Inte kontrollerade; se \cref{sec:fortsatt} \\
    \bottomrule
  \end{tabular}
\end{table}

Andra raden förtjänar en kommentar.  Påståendena om lärande i marginalens
lärarnoter vilar på Martons \emph{Necessary Conditions of Learning}.
Citaten på sidorna 45--47 är kontrollerade mot boken; citatet på sidan 89
är hämtat ur en andrahandsuppteckning som anger boken och sidan, och är
alltså inte läst i original här.  Det står i litteraturlistans källfil,
och det står också i tabellen --- för det som inte är kontrollerat ska
räknas upp lika noga som det som är det.

\section{Gör påståendet till en fråga}

Ett påstående går att kontrollera först när det är tydligt vad som skulle
kunna visa att det är \emph{fel}.  Formulera det därför som en fråga med
ett svar: \enquote{svarar nybörjare 5 på frågan om det största indexet i
en lista med fem element?} kan besvaras med ja eller nej, medan
\enquote{listindex är svårt} inte kan det.  Ingen tänkbar källa skulle
visa att det senare är fel, eftersom \enquote{svårt} är ett omdöme och
inte ett faktum.

Ofta döljer sig flera påståenden i ett.  \enquote{Principen DRY är bra och
kommer från Hunt och Thomas} är två frågor --- \emph{håller den?} och
\emph{varifrån kommer den?} --- och de kräver olika slags källor.

Det gäller särskilt \emph{råd}.  Varje gång ett material säger att du ska
göra något på ett visst sätt påstår det samtidigt att det blir bättre så,
och det är ett påstående som kan prövas.  De fem råden i
\cref{sec:fortsatt} är alla av det slaget.

\section{Välj rätt sorts källa}

Olika frågor besvaras av olika källor.
\begin{description}
  \item[Uppslagsverk] för vad ett ord \emph{betyder}.
  \item[Primärkällan] för var en idé \emph{kommer ifrån}.  Vill man veta
    om DRY härrör från Hunt och Thomas läser man deras bok, inte en
    lärobok som återger den.
  \item[Normativ dokumentation] för vad ett verktyg \emph{gör}.  Att
    \mintinline{python}{random.shuffle} blandar en lista på plats är
    ingen forskningsfråga; det står i Pythons dokumentation, och det går
    dessutom att köra.
  \item[Översikter] (\foreignlanguage{english}{reviews}) för om något är
    \emph{etablerat}.  En systematisk översikt går igenom hundratals
    studier efter en angiven metod och väger därför tyngre än en enskild
    studie.
  \item[Enskilda studier] för ett specifikt resultat --- med förbehållet
    att en enskild studie bara visar att något observerats en gång.
\end{description}
Var man hittar dem: universitetsbibliotekets söktjänst, och databaser som
Scopus, Web of Science, IEEE Xplore, ACM Digital Library, DBLP (datalogi)
och OpenAlex (OA; öppen, och inte samma sak som
\foreignlanguage{english}{open access}).  Google Scholar täcker brett men
rankar ogenomskinligt.

\section{Sök så att du kan ha fel}

Den som söker på namn hon redan känner till hittar det hon väntar sig.
Sök därför \emph{ämnesdrivet} --- på begrepp, inte författare --- och
ställ samma frågor i alla databaser.

Sök också efter \emph{invändningar}: lägg till ord som
\enquote{critique}, \enquote{limitations} eller \enquote{failure}.  Ett
påstående som överlevt en motsökning är värt mer än ett som aldrig
prövats, och en invändning som hittas blir ett \emph{förbehåll} i svaret,
inte ett nederlag.

\section{Läs i fulltext, och skriv ner hur}

En träff i en databas är inte ett belägg.  Sammanfattningen
(\foreignlanguage{english}{abstract}) räcker inte heller: den säger vad
författarna vill ha sagt, inte vad de visat.  Läs hela texten och leta
upp \emph{det ställe} som styrker påståendet --- ett ordagrant citat med
sidnummer.  Räkna aldrig en källa du inte läst.

Det gäller också en källa du ärvt, och det är just den fällan den här
övningen står närmast.  Att ett påstående står med en fotnot i ett annat
material betyder inte att du har kontrollerat det; du har kontrollerat
att någon annan påstod det.  Övningen återanvänder därför bara sådana
påståenden som är kontrollerade \emph{i fulltext} i en annan bilaga, och
pekar ut vilken --- och där ett citat kommer från andra hand, som på sidan
89 hos Marton, står det utskrivet.

Skriv sedan ner hur du gick till väga.  I det här materialet står den
anteckningen som en kommentar ovanför varje post i litteraturlistans
källfil, med fasta fält:
\begin{description}
  \item[\texttt{CLAIM}] vilket påstående källan ska styrka;
  \item[\texttt{FOUND-VIA}] hur den hittades --- sökfråga och databas,
    eller vilken bilaga den är kopierad från;
  \item[\texttt{PICKED}] varför just den, framför andra träffar;
  \item[\texttt{QUOTE}] det ordagranna citatet, med sida;
  \item[\texttt{VERIFIED}] att fulltexten lästs, och var;
  \item[\texttt{COUNTER}] vad motsökningen gav;
  \item[\texttt{DATE}] när.
\end{description}
Det tar ett par minuter per källa, och det är det enda som gör det
möjligt att om ett år svara på frågan \enquote{hur vet du det?}

\section{Dra slutsatsen --- med förbehåll}

Svaret skrivs ut tillsammans med det som talar emot.  Förbehållen är inte
svagheter i svaret; de \emph{är} svaret, och de avgör hur påståendet får
användas.  Ett påstående vars belägg bara räcker till
sammanfattningsnivå ska formuleras svagare än ett som är läst i fulltext,
och det ska synas att så är fallet.  Det är därför inledningen till
övningen säger att vinsten av att försöka själv först är \emph{måttlig}
och inte \emph{stor}.

\section{En anmärkning om språkmodeller}

En språkmodell är användbar för att \emph{hitta} och \emph{sortera}
källor, och oanvändbar som sista led.  Den kan uppge en sida som inte
finns, eller ett citat som inte står där.  Varje källa som citeras ska
läsas och väljas av en människa, och den människan är den som får svara
för innehållet.

\section{Fortsatt arbete}
\label{sec:fortsatt}

Sex påståenden i övningen är inte kontrollerade, och de är av två slag.

Fem av dem är \emph{råd}: att en namngiven konstant är lättare att ändra
än ett tal inne i koden och att ett namn som säger vad funktionen lämnar
tillbaka gör programmet lättare att läsa (\cref{sec:fem-fel}); att en
\mintinline{python}{for}-slinga över en lista är lättare att läsa och
svårare att skriva fel än en \mintinline{python}{while}-slinga med en
räknare (\cref{sec:listan}); att det är en färdighet i sig att hitta och
läsa dokumentationen (\cref{sec:blandat}); och att en nästlad slinga blir
lättare att läsa när den inre slingan flyttas till en egen funktion
(\cref{sec:multtabell}).  Alla fem påstår en nytta, och veckans
föreläsningar gör flera av dem utan belägg.  Frågan att söka på är i
varje fall densamma: finns det någon studie som har mätt den nyttan, och
vad säger de som har försökt och inte funnit den?

Det sjätte är en \emph{missuppfattning} som namnges utan källa: att 5 är
det vanliga felsvaret på frågan om det största indexet i en lista med fem
element (\cref{sec:listan}).  Föreläsningen \emph{Behållare: Listor}
citerar två källor om listindex, men har ingen granskad bilaga att peka
på, så påståendet kan inte återanvändas härifrån.

Därtill återstår tre saker att göra.  Aritmetikens fundamentalsats anges i
\cref{sec:primtal} utan källa, vilket är enkelt att rätta men ändå inte
gjort.  Citatet från sidan 89 hos Marton bör läsas i original, så att
andrahandsuppteckningen kan strykas ur källfilen.  Och samma citat
används med samma brist i föreläsningarna \emph{Behållare: Listor} och
\emph{Behållare: Tupler}, så rättelsen hör hemma där också.
